\begin{afterword}
\summaryhead

The author imagines performing real sorcery, a light conjured from empty space, and keeping everyone from being curious about it by answering every question with ``Science!'' The word explains nothing: it does not say whether the light will brighten or fade, or how to make one. Yet people turn away satisfied. ``Electricity!'' does the same for an ordinary light bulb.

Being told that something is ``scientifically explicable,'' the post argues, means only that someone else knows how it works, and this adds nothing to what you know. It is not greed that makes the difference, since the post is concerned only with the feeling of curiosity. Why should your curiosity shrink because someone else knows? If it did, a voice saying ``Somebody already knows why this elephant is here'' would let you walk past a hovering green elephant. The author admits not knowing where the kidneys are, and says most readers do not know how their cerebellum moves their hand. The world is full of puzzles you personally do not understand, so you should not complain that science has emptied it of mystery.

\responsehead

The post's starting point is sound and continues the sequence. A word that names a field of knowledge, ``Science!'' or ``Electricity!'', gives the listener no expectations about the thing in front of them. The sequence's earlier posts had called such answers curiosity-stoppers and passwords, and this post extends the idea to answers that are true.

The central step leaves something out. The post takes ``scientifically explicable'' to mean that someone else knows, and says that the listener then knows nothing more. About how the bulb will behave, that is right. But the phrase also tells the listener that the bulb breaks no known law, and that an answer can be found. A light that science could not explain would show that something the listener took as settled is mistaken. That is a reason for it to command more attention than an explained one, and it has nothing to do with who else knows.

The post does not consider this reason. It answers the objection in the form of fame and fortune, and then asks whether what is left is spite. The question comes with no evidence. The post's own examples fit the reason it leaves out. The green elephant holds attention because it cannot be fitted into anything the reader knows, whatever the voice says. The light spell would intrigue more than a raised hand because the spell would break known physics and the hand does not. The post calls that difference a ``double standard''; it is the difference between an event that fits what you know and one that does not.

The conclusion is narrower than the complaint it answers. It tells the reader that much remains unknown to them, which says nothing about the things they do come to understand. The author gave that fuller answer later, in ``Joy in the Merely Real.''

In short: a sound point about names that pass for explanations, together with the suggestion that curiosity which drops when someone else knows is spite, when the plainer reason is that an explained event tells you that nothing you believed is mistaken.
\end{afterword}
